\documentclass[12pt,a4paper]{report}

% --- Pacchetti ---
\usepackage[utf8]{inputenc}
\usepackage[T1]{fontenc}
\usepackage[a4paper,margin=2.8cm]{geometry}

% Etichette in italiano (babel-italian non disponibile in questo ambiente)
\renewcommand{\contentsname}{Indice}
\renewcommand{\chaptername}{Capitolo}
\renewcommand{\bibname}{Bibliografia}
\renewcommand{\tablename}{Tabella}
\renewcommand{\figurename}{Figura}
\usepackage{graphicx}      % Inserimento immagini
\usepackage{booktabs}      % Tabelle di qualità
\usepackage{longtable}     % Tabelle su più pagine
\usepackage{array}
\usepackage{enumitem}
\usepackage{fancyhdr}
\usepackage[hidelinks]{hyperref}

% --- Intestazioni ---
\pagestyle{fancy}
\fancyhf{}
\fancyhead[L]{\leftmark}
\fancyhead[R]{\thepage}
\renewcommand{\headrulewidth}{0.4pt}

% --- Metadati ---
\title{Progettazione di una ECU \emph{software-defined} per motori
monocilindrici:\\ architettura, prototipo e validazione}
\author{}
\date{Giugno 2026}

\begin{document}

\maketitle

% ============================================================
\chapter*{Sommario}
\addcontentsline{toc}{chapter}{Sommario}
\markboth{Sommario}{Sommario}

Questo elaborato presenta la progettazione di una centralina elettronica
(\emph{Electronic Control Unit}, ECU) \emph{custom} per motori monocilindrici,
concepita secondo un approccio \emph{software-defined}. Il lavoro non si propone
come prodotto finito, ma come \textbf{prima iterazione} e \emph{blueprint}: un
percorso che parte dalla definizione dell'idea progettuale, prosegue con
l'analisi dei requisiti, la progettazione di un'architettura indipendente dalla
specifica realizzazione, lo sviluppo di un prototipo focalizzato su un
sottoinsieme significativo di funzionalità, e si chiude con una fase di
validazione mirata a verificare la coerenza tra prototipo e requisiti.

L'obiettivo non è dimostrare prestazioni motoristiche, bensì mostrare che
un'architettura modulare, testabile e aperta all'evoluzione costituisce una base
solida per la gestione di un motore monocilindrico e per i servizi di
telemetria, tuning e analisi storica che vi gravitano attorno. Le scelte di
progetto sono inquadrate rispetto allo stato dell'arte delle centraline
programmabili, dei sistemi embedded real-time e dei \emph{digital twin}
motoristici.

\tableofcontents

% ============================================================
\chapter{Contesto e idea di progetto}
\label{ch:contesto}

\section{Il dominio applicativo}
La gestione elettronica di un motore a combustione interna consiste nel
trasformare un insieme di segnali sensoriali, che descrivono lo stato
istantaneo del propulsore e l'intenzione del pilota, in una sequenza di comandi
attuativi temporalmente precisi. Nei motori monocilindrici a carburatore di
piccola e media cilindrata, tipici di motociclette, kart e applicazioni
sportive amatoriali, la grandezza di controllo principale è l'\emph{anticipo di
accensione}, ovvero l'istante in cui viene comandata la scintilla rispetto alla
posizione angolare dell'albero motore. A questa si affiancano funzioni
accessorie quali l'arricchimento della miscela agli alti regimi, la gestione di
una eventuale valvola di scarico e strategie di sicurezza orientate alla
protezione termica del motore.

A differenza dei moderni propulsori automobilistici, in questo segmento la
centralina di serie è spesso un dispositivo a comportamento fisso, non
riprogrammabile e privo di interfacce di diagnosi. Chi interviene sul motore per
finalità di messa a punto si trova quindi a operare con strumenti limitati,
oppure deve ricorrere a centraline aftermarket dal costo e dalla complessità non
sempre giustificati per applicazioni di piccola scala.

\section{L'idea di progetto}
L'idea alla base di questo lavoro è la realizzazione di una ECU \emph{custom},
\emph{software-defined}, destinata a motori monocilindrici a carburatore (sia 2T
che 4T). Il termine \emph{software-defined} indica che il comportamento del
sistema non è cablato nell'hardware, ma definito ed evolvibile via software: le
strategie di controllo, le mappe di accensione e i parametri di sicurezza
risiedono in dati e logica modificabili, aggiornabili anche da remoto.

Il valore dell'idea non sta nel \emph{come} viene controllata la singola
scintilla, ma nel \emph{cosa} il sistema rende possibile:
\begin{itemize}[noitemsep]
  \item un controllo motore deterministico, separato in modo netto dai
        compiti non critici;
  \item un canale di telemetria in tempo reale verso un'interfaccia di tuning
        accessibile da browser, senza software dedicato;
  \item la capacità di modificare le mappe e i parametri ``a caldo'' e di
        valutarne gli effetti;
  \item la registrazione delle sessioni di funzionamento e la loro analisi
        storica, fino alla costruzione di una rappresentazione persistente del
        comportamento della centralina (\emph{digital twin});
  \item l'aggiornamento del firmware \emph{over-the-air} (OTA).
\end{itemize}

In altre parole, il progetto sposta il baricentro dalla centralina intesa come
componente elettronico statico verso un \emph{sistema} il cui valore cresce nel
tempo attraverso il software e i dati, in analogia con la tendenza più ampia dei
veicoli \emph{software-defined}.

\section{Inquadramento rispetto allo stato dell'arte}
L'idea trova riscontro in tre filoni, accademici e industriali, che ne
confermano la pertinenza.

\paragraph{Centraline programmabili e aftermarket.}
Le ECU programmabili costituiscono una categoria consolidata: a differenza delle
centraline a comportamento fisso, consentono all'utente di definire mappe di
accensione e di combustibile su misura del motore, tipicamente tramite un
calcolatore collegato durante il funzionamento~\cite{linkecu,patent_prog}. Il
presente lavoro si colloca in questo solco, ma con due differenze qualificanti:
l'apertura dell'architettura software e l'integrazione nativa di telemetria,
analisi storica e aggiornamento remoto, funzioni generalmente assenti o
proprietarie nelle soluzioni di fascia bassa.

\paragraph{Sistemi embedded real-time multicore.}
Il controllo motore richiede determinismo temporale (\emph{hard real-time}) per
la gestione degli eventi di accensione. L'adozione di un sistema operativo
real-time consente di schedulare e prioritizzare i task garantendo la
periodicità del ciclo di controllo, e le architetture multicore permettono di
isolare i compiti critici da quelli di comunicazione~\cite{freertos_idf,
freertos_perf}. Questo conferma la praticabilità di una separazione netta tra un
nucleo di controllo deterministico e un sottosistema di connettività, principio
cardine dell'architettura proposta.

\paragraph{Validazione \emph{hardware-in-the-loop}.}
La validazione delle centraline automobilistiche si avvale ormai
sistematicamente del \emph{hardware-in-the-loop} (HIL), in cui l'hardware reale
esegue il firmware di produzione mentre interagisce con un modello simulato
dell'impianto fisico; ciò consente di rendere testabili tempistiche, carichi e
guasti prima dell'integrazione sul veicolo~\cite{hil_mathworks, hil_framework}.
La presenza, nel progetto, di un nodo simulatore e di un \emph{harness}
sensoriale commutabile tra sorgenti reali e sintetiche riprende, in scala
ridotta, questo paradigma.

\paragraph{Digital twin motoristici.}
Infine, la letteratura recente documenta l'impiego di \emph{digital twin} dei
motori a combustione interna per il monitoraggio, la diagnostica e la modellazione
predittiva, costruiti a partire da dati sensoriali trasmessi a un
server~\cite{dt_review, dt_diesel}. L'obiettivo di lungo periodo del progetto --- una
rappresentazione persistente lato server che ricostruisce ciò che la ECU ha
osservato, deciso e comandato --- è coerente con questa direzione, pur partendo
da una versione descrittiva e orientata al \emph{replay}.

\section{Natura del lavoro: una prima iterazione}
È importante chiarire l'orizzonte dell'elaborato. Non si tratta di una centralina
pronta per l'impiego su strada o in pista, né di un sistema validato su motore
reale. Si tratta di un \textbf{primo prototipo} e di una \emph{base
architetturale}: alcune scelte sono deliberatamente semplificate (per esempio
l'uso di mappe monodimensionali e la simulazione di alcuni sensori), e parte
delle funzioni è rinviata a sviluppi futuri. Il contributo risiede nella
definizione di un'architettura coerente e testabile e nella dimostrazione, su un
sottoinsieme di funzionalità, che tale architettura regge alla prova della
realizzazione.

% ============================================================
\chapter{Analisi dei requisiti}
\label{ch:requisiti}

L'analisi dei requisiti distingue \emph{cosa} il sistema deve fare (requisiti
funzionali) da \emph{come} deve farlo, in termini di qualità (requisiti non
funzionali). I requisiti sono espressi a livello di sistema, indipendentemente
dalla specifica realizzazione del prototipo.

\section{Attori e casi d'uso principali}
Il sistema coinvolge tre attori:
\begin{itemize}[noitemsep]
  \item il \textbf{motore}, sorgente dei segnali fisici e destinatario dei
        comandi attuativi;
  \item l'\textbf{utente} (pilota o tecnico di messa a punto), che osserva lo
        stato, modifica le mappe e impartisce comandi;
  \item il \textbf{server}, che conserva la storia operativa e ne abilita
        l'analisi.
\end{itemize}
I casi d'uso principali sono: controllo dell'accensione in tempo reale,
monitoraggio della telemetria, modifica e selezione delle mappe, registrazione e
revisione delle sessioni, aggiornamento del firmware.

\section{Requisiti funzionali}
\begin{longtable}{p{1.6cm} p{4.2cm} p{7.5cm}}
\toprule
\textbf{ID} & \textbf{Requisito} & \textbf{Descrizione} \\
\midrule
\endhead
RF-1 & Sincronizzazione motore & Il sistema deve ricostruire posizione e velocità
angolare del motore a partire dal segnale di pick-up, riconoscendo l'avvenuta
sincronizzazione. \\
RF-2 & Calcolo dell'anticipo & Il sistema deve determinare l'anticipo di
accensione in funzione del punto operativo, mediante interpolazione su mappa. \\
RF-3 & Comando di accensione & Il sistema deve comandare l'evento di accensione
(CDI) con la corretta temporizzazione rispetto al riferimento di fase. \\
RF-4 & Acquisizione ingressi & Il sistema deve acquisire posizione farfalla
(TPS), temperatura gas di scarico (EGT) e ulteriori ingressi (quick shifter,
selettore mappa). \\
RF-5 & Attuazione ausiliaria & Il sistema deve comandare in PWM gli attuatori
ausiliari (Power Jet ed eventuale valvola di scarico) in funzione del punto
operativo. \\
RF-6 & Strategie di sicurezza & Il sistema deve applicare protezioni (allarme
EGT, limitatore di giri, taglio accensione per quick shifter). \\
RF-7 & Gestione delle mappe & Il sistema deve consentire creazione, modifica,
selezione e persistenza di più mappe. \\
RF-8 & Telemetria in tempo reale & Il sistema deve trasmettere a un client lo
stato corrente (giri, TPS, EGT, stato, anticipo, duty) ad alta frequenza. \\
RF-9 & Comandi da remoto & Il sistema deve ricevere comandi dal client (cambio
mappa, modifica mappa, trigger quick shifter, richiesta configurazione). \\
RF-10 & Registrazione sessioni & Il sistema deve registrare le sessioni di
funzionamento e renderle disponibili per l'analisi storica. \\
RF-11 & Aggiornamento OTA & Il sistema deve poter verificare e applicare
aggiornamenti del firmware da remoto. \\
RF-12 & Interfaccia utente & Il sistema deve offrire un'interfaccia di tuning e
monitoraggio accessibile senza software dedicato. \\
\bottomrule
\end{longtable}

\section{Requisiti non funzionali}
\begin{longtable}{p{1.6cm} p{4.2cm} p{7.5cm}}
\toprule
\textbf{ID} & \textbf{Qualità} & \textbf{Descrizione} \\
\midrule
\endhead
RNF-1 & Determinismo & Le funzioni di controllo motore devono essere
\emph{hard real-time}; la loro temporizzazione non deve dipendere dai compiti di
comunicazione. \\
RNF-2 & Separazione delle responsabilità & I compiti critici e non critici
devono essere isolati, idealmente su unità di elaborazione distinte. \\
RNF-3 & Indipendenza del nucleo & Il controllo motore deve restare pienamente
operativo anche in assenza di client, server e connettività. \\
RNF-4 & Modularità e testabilità & La logica di dominio deve essere indipendente
dall'hardware, così da poter essere verificata su calcolatore senza il
dispositivo fisico. \\
RNF-5 & Evolvibilità & L'architettura deve permettere di aggiungere sensori,
attuatori e strategie senza riprogettare il sistema. \\
RNF-6 & Vincoli di risorsa & Il sistema deve operare entro i limiti di memoria e
flash della piattaforma scelta. \\
RNF-7 & Tracciabilità & Ogni dato storico deve essere riconducibile alla
versione di firmware, mappe e configurazione in uso al momento della
registrazione. \\
RNF-8 & Sicurezza dei dati & Le operazioni di replay e simulazione non devono
mai sovrascrivere i dati originali registrati. \\
RNF-9 & Accessibilità & L'interfaccia utente deve essere fruibile da dispositivi
comuni (telefono, tablet) tramite browser. \\
\bottomrule
\end{longtable}

\section{Confini di scopo}
Coerentemente con la natura di prima iterazione, alcune funzioni rientrano nello
scopo del prototipo e altre sono esplicitamente rinviate. Restano fuori scopo, e
costituiscono sviluppi futuri: il sensore di detonazione con correzione attiva,
il controllo della valvola di scarico, le mappe tridimensionali
(giri~$\times$~farfalla), la pipeline completa di integrazione continua HITL, il
progetto della scheda dedicata (PCB) e l'integrazione di sensori reali oggi
simulati. Il dettaglio del confine di scopo è ripreso nel
Capitolo~\ref{ch:conclusioni}.

% ============================================================
\chapter{Progettazione dell'architettura}
\label{ch:architettura}

Questo capitolo descrive l'architettura del sistema in termini di macro
sottosistemi e di relazioni tra essi. La descrizione è volutamente indipendente
dalla specifica realizzazione: prescinde dal prototipo e mira a catturare la
struttura logica che ne resta valida anche al variare delle tecnologie.

\section{Decomposizione in macro sottosistemi}
Il sistema si articola in quattro macro sottosistemi, organizzati su un confine
netto tra ciò che è critico per il controllo motore e ciò che non lo è.

\begin{description}
  \item[Nodo di controllo (ECU).] È l'autorità sullo stato real-time del motore.
        Acquisisce i sensori, ricostruisce sincronizzazione e velocità, esegue
        le strategie di controllo, schedula gli eventi di accensione e di
        attuazione, applica le protezioni e produce telemetria. Deve poter
        funzionare in completa autonomia.
  \item[Sottosistema di simulazione (banco prova).] Genera segnali sensoriali
        sintetici per sollecitare la ECU lungo il normale percorso di ingresso,
        consentendo prove ripetibili in assenza del motore reale.
  \item[Client (interfaccia utente).] Presenta lo stato corrente, consente
        l'editing e la selezione delle mappe, invia comandi alla ECU e funge da
        ponte di comunicazione verso il server. Non è il proprietario permanente
        dei dati storici.
  \item[Server e \emph{digital twin}.] Mantiene la rappresentazione persistente
        della ECU: identità, revisioni di firmware/mappe/configurazione, sessioni
        registrate, e i servizi di analisi, \emph{replay} e simulazione.
\end{description}

\section{Comunicazione tra sottosistemi}
Le relazioni tra i sottosistemi seguono direzioni e responsabilità ben definite,
riassunte in Tabella~\ref{tab:comunicazioni}. Un punto architetturale rilevante è
che la ECU non dispone di connettività diretta verso il server: la comunicazione
remota transita attraverso il client, che agisce da \emph{gateway}. Ciò mantiene
il nodo di controllo semplice e indipendente da Internet.

\begin{table}[h]
\centering
\caption{Relazioni di comunicazione tra i macro sottosistemi.}
\label{tab:comunicazioni}
\small
\begin{tabular}{p{4.2cm} p{3.2cm} p{5.2cm}}
\toprule
\textbf{Percorso} & \textbf{Natura} & \textbf{Contenuto} \\
\midrule
Simulatore $\rightarrow$ ECU & Segnali fisici & Onda quadra (pick-up), tensioni
analogiche (TPS, EGT) \\
ECU $\leftrightarrow$ Client & Canale bidirezionale a bassa latenza & Telemetria
in push; comandi e modifiche mappa in pull \\
Client $\rightarrow$ Server & Canale Internet & Inoltro di sessioni e telemetria
rilevante \\
ECU $\leftrightarrow$ Server & Mediata dal client & Verifica e download
aggiornamenti firmware \\
\bottomrule
\end{tabular}
\end{table}

\section{Architettura interna del nodo di controllo}
Per soddisfare i requisiti di determinismo (RNF-1), separazione (RNF-2) e
testabilità (RNF-4), il nodo di controllo è progettato secondo il principio di
\emph{separazione tra dominio e infrastruttura}, riconducibile al modello a
porte e adattatori (\emph{ports and adapters}). Le dipendenze puntano verso
l'interno: il nucleo di controllo non conosce le tecnologie di comunicazione o
di memorizzazione, ma soltanto interfacce astratte.

\begin{description}
  \item[Dominio.] Tipi di valore, macchine a stati, algoritmi e politiche, scritti
        senza alcun riferimento all'hardware. Vi risiedono la macchina a stati del
        motore, la stima della velocità, le mappe con interpolazione e i criteri
        di sicurezza.
  \item[Applicazione.] Orchestrazione dei casi d'uso (controllo accensione,
        attuazione ausiliaria, gestione mappe, registrazione e pubblicazione delle
        sessioni, instradamento dei comandi).
  \item[Porte.] Interfacce astratte che definiscono i confini del sistema
        (ingressi sensoriali, schedulazione dell'accensione, persistenza delle
        mappe, trasporto della telemetria e delle sessioni).
  \item[Infrastruttura.] Realizzazioni concrete delle porte sulla piattaforma
        scelta.
  \item[Runtime.] Creazione dei task, assegnazione alle unità di elaborazione e
        cablaggio tra porte e adattatori.
\end{description}

Questa stratificazione è ciò che rende il sistema testabile su calcolatore: il
dominio, non dipendendo dall'hardware, può essere esercitato con sorgenti
sostitutive e verificato in modo automatico, come illustrato nel
Capitolo~\ref{ch:validazione}.

\section{Modello dello stato del motore}
Il comportamento del motore è descritto da una macchina a stati finiti con sei
stati: inizializzazione, sincronizzazione, funzionamento, minimo, taglio
accensione e allarme. Le transizioni sono guidate da eventi (primo impulso
rilevato, sincronizzazione acquisita o persa, superamento di soglie, richiesta
di quick shifter). La macchina a stati è il punto in cui convergono controllo e
sicurezza: per esempio, il superamento della soglia EGT porta allo stato di
allarme, mentre una richiesta di quick shifter induce un taglio temporaneo
dell'accensione.

\section{Flusso dei dati di controllo}
Il flusso logico del controllo, indipendente dalla realizzazione, è il seguente:
il rilevamento di un impulso di pick-up innesca l'acquisizione rapida di
velocità e ingressi analogici; segue la determinazione dell'anticipo per
interpolazione sulla mappa attiva, con le eventuali correzioni di sicurezza; il
risultato è tradotto in una temporizzazione dell'evento di accensione e nei duty
cycle degli attuatori ausiliari; infine, lo stato corrente è reso disponibile,
attraverso un'area dati condivisa, al sottosistema di comunicazione che lo
trasmette al client. Il principio chiave è che l'esecuzione degli eventi
critici non dipende dai compiti non critici.

\section{Il digital twin come orizzonte architetturale}
Sul lato server, l'architettura prevede una rappresentazione persistente della
ECU --- il \emph{digital twin} --- che combina identità e revisioni del
dispositivo con la storia operativa registrata. Concettualmente, il twin
distingue la \emph{provenienza} di ogni valore: misurato, derivato, deciso,
comandato, stimato o simulato. Questa distinzione abilita tre operazioni
mantenute separate: la riproduzione dei dati registrati, il \emph{replay}
deterministico della logica di controllo sugli stessi ingressi, e la simulazione
di mappe o parametri alternativi. In tutti i casi, i dati originali non vengono
mai sovrascritti (RNF-8) e ogni risultato resta tracciabile alla configurazione
che lo ha prodotto (RNF-7). Nel prototipo questo orizzonte è realizzato solo in
forma descrittiva e orientata al \emph{replay}, lasciando i modelli avanzati agli
sviluppi futuri.

% ============================================================
\chapter{Prototipo e scelte implementative}
\label{ch:prototipo}

Il prototipo realizza un \emph{sottoinsieme} significativo dell'architettura,
sufficiente a dimostrarne la praticabilità. Questo capitolo motiva le scelte
principali, senza scendere in dettagli implementativi non necessari alla
comprensione.

\section{Piattaforma di calcolo}
Il nodo di controllo è basato sul SoC ESP32-S3, montato su una scheda di sviluppo
compatta. La scelta è motivata dal bilanciamento tra capacità di calcolo (due
core), connettività nativa (Wi-Fi e Bluetooth LE), disponibilità di periferiche
temporizzate e costo contenuto. Le caratteristiche salienti della scheda sono
riportate in Tabella~\ref{tab:board}.

\begin{table}[h]
\centering
\caption{Caratteristiche principali della piattaforma di prototipazione.}
\label{tab:board}
\small
\begin{tabular}{p{4.5cm} p{8.5cm}}
\toprule
\textbf{Elemento} & \textbf{Valore} \\
\midrule
SoC & ESP32-S3 (Xtensa LX7 dual-core, fino a 240 MHz) \\
Connettività & Wi-Fi 2.4 GHz, Bluetooth LE \\
Flash / PSRAM & 4 MB / 2 MB \\
Tensione logica & 3,3 V \\
GPIO esposti & 27 pin oltre alle linee di alimentazione \\
\bottomrule
\end{tabular}
\end{table}

\subsection{Firmware nativo e sistema real-time}
Il firmware è sviluppato in ambiente nativo ESP-IDF su FreeRTOS, senza ricorrere
al framework Arduino. La motivazione è il controllo fine sulla schedulazione e il
determinismo: FreeRTOS, integrato in ESP-IDF, consente di creare task con
priorità e di assegnarli a uno specifico core. I compiti critici di controllo
sono collocati su un core, quelli di comunicazione sull'altro, realizzando la
separazione richiesta dai requisiti RNF-1 e RNF-2.

\section{Sottosistema sensoriale e \emph{harness} commutabile}
Il prototipo adotta un'astrazione delle sorgenti sensoriali in cui ogni sensore
è servito attraverso un servizio di dominio indipendente dall'hardware. Un
\emph{harness} permette di selezionare, in fase di configurazione, se ciascun
ingresso provenga da una sorgente reale o da una sorgente \emph{fake}
deterministica. In modalità predefinita tutte le sorgenti sono sintetiche e
producono campioni riproducibili, stampati in formato CSV su seriale per
ispezione e tracciamento.

Questa scelta è centrale: rende la stessa logica esercitabile sia su banco sia
con segnali reali, e costituisce il fondamento della strategia di validazione.
Allo stato attuale del prototipo, gli ingressi confermati come reali sono la
posizione farfalla (via ADC), il quick shifter, il selettore mappa e il segnale
di pick-up (via cattura su periferica temporizzata), mentre temperatura acqua,
EGT e detonazione restano simulati in attesa del relativo cablaggio.

\section{Strategie di controllo realizzate}
Il prototipo realizza la macchina a stati del motore, la stima della velocità a
partire dagli istanti del pick-up, il calcolo dell'anticipo e del duty del Power
Jet mediante mappe \emph{monodimensionali} con interpolazione lineare, e le
strategie di sicurezza essenziali (allarme EGT, taglio accensione per quick
shifter). La scelta di mappe 1D, anziché 3D, è una semplificazione deliberata
adeguata alla prima iterazione: riduce la complessità di editing e di
persistenza, pur mantenendo invariata la struttura dei moduli che le utilizzano.
Le mappe sono persistite nella memoria non volatile e ne è possibile la
selezione e la modifica a runtime.

\section{Connettività e interfaccia utente}
Il nodo di controllo espone un server che gestisce, sul medesimo punto di
accesso, la distribuzione dei contenuti statici dell'interfaccia e un canale
bidirezionale a bassa latenza per la telemetria e i comandi. L'interfaccia di
tuning è un'applicazione web servita direttamente dalla ECU, accessibile da
telefono o tablet senza installazione di software dedicato. Per contenere
l'occupazione di flash, gli asset statici sono memorizzati compressi e serviti
con la relativa codifica, lasciando al browser la decompressione trasparente.
La telemetria viaggia in \emph{push} ad alta frequenza, mentre i comandi (cambio
e modifica mappa, trigger quick shifter, richiesta configurazione) viaggiano
nella direzione opposta.

\section{Registrazione delle sessioni e budget di memoria}
La registrazione delle sessioni è vincolata dalla memoria disponibile sul
dispositivo. Ogni campione è codificato in forma binaria compatta e accumulato
in un buffer; la frequenza di campionamento è scelta come compromesso tra
fluidità dei grafici e durata massima della sessione. Al termine della sessione,
i dati sono trasmessi verso il server attraverso il client, eventualmente
suddivisi in blocchi. Si distingue tra la telemetria ad alta frequenza, destinata
alla visualizzazione in tempo reale e non registrata, e i campioni di sessione a
frequenza inferiore, destinati all'analisi storica.

\section{Nodo simulatore}
Il banco prova è realizzato con un secondo microcontrollore che genera i segnali
attesi dalla ECU: un'onda quadra a frequenza regolabile per il pick-up e
tensioni analogiche per TPS ed EGT. Una piccola interfaccia web dedicata
consente di impostare i valori e di richiamare condizioni operative predefinite
(minimo, crociera, pieno carico) oltre a un comando di arresto. Il simulatore
sollecita la ECU lungo lo stesso percorso di ingresso dei sensori reali, in
coerenza con l'approccio HIL discusso nel Capitolo~\ref{ch:contesto}.

% ============================================================
\chapter{Validazione}
\label{ch:validazione}

La validazione verifica che il prototipo soddisfi i requisiti funzionali
individuati nel Capitolo~\ref{ch:requisiti}. La strategia è strutturata su più
livelli, dal più astratto e veloce al più vicino all'hardware, sfruttando
l'astrazione delle sorgenti sensoriali e la separazione tra dominio e
infrastruttura.

\section{Strategia complessiva}
La testabilità è un requisito non funzionale di primo piano (RNF-4), e
l'architettura è stata progettata per soddisfarlo. La validazione segue tre
livelli complementari:
\begin{enumerate}[noitemsep]
  \item test su calcolatore (\emph{host}) della logica di dominio, senza
        hardware;
  \item esecuzione sul dispositivo con sorgenti interamente sintetiche
        (\emph{all-fake});
  \item esecuzione sul dispositivo con sorgenti miste, reali e sintetiche, fino
        alla verifica dei singoli ingressi fisici.
\end{enumerate}

\section{Test su calcolatore}
Il primo livello esercita il dominio sensoriale, i contratti dei servizi, le
sorgenti sintetiche, l'archivio dati e la formattazione dell'output, senza
richiedere né ESP-IDF né hardware. È il livello più rapido e ripetibile, ed è
quello che rende sostenibile l'iterazione sul software. I test di livello driver
per i sensori su bus SPI sono volutamente rinviati e la loro logica è verificata
indirettamente tramite campioni iniettati a livello di dominio. Questo livello
dà evidenza diretta del soddisfacimento dei requisiti di calcolo: stima della
velocità (RF-1), interpolazione delle mappe (RF-2), gestione delle mappe (RF-7) e
logica di sicurezza (RF-6).

\section{Esecuzione con sorgenti sintetiche}
Il secondo livello esegue il firmware completo sul dispositivo con tutte le
sorgenti sintetiche. Iniettando campioni deterministici di tutti gli ingressi, si
osserva la risposta integrata del sistema: la velocità diventa valida dopo un
numero sufficiente di catture, gli ingressi analogici variano nel tempo e gli
eventi discreti (quick shifter, cambio mappa) compaiono nei punti attesi. Questo
livello valida l'acquisizione integrata degli ingressi (RF-4) e il
comportamento della macchina a stati lungo le transizioni nominali.

\section{Esecuzione con sorgenti miste}
Il terzo livello commuta selettivamente singoli ingressi da sintetici a reali,
verificando un sensore alla volta. Per esempio, una rampa di tensione
sull'ingresso farfalla deve riflettersi nella lettura corrispondente; la
commutazione del quick shifter o del selettore mappa deve generare i relativi
eventi; un'onda quadra applicata all'ingresso di pick-up a frequenza nota deve
produrre il valore di velocità atteso, e la perdita di un impulso deve rendere
``stale'' la lettura prima del recupero. Questo livello collega la logica già
validata ai segnali fisici reali, dando evidenza del percorso completo di
sincronizzazione (RF-1) e acquisizione (RF-4).

\section{Verifiche di confine}
A garanzia dei requisiti di modularità ed evolvibilità (RNF-4, RNF-5), sono
previste verifiche automatiche che controllano che il dominio sensoriale resti
indipendente dai driver di piattaforma e dai sottosistemi non sensoriali (per
esempio telemetria, persistenza, attuazione). L'assenza di dipendenze indebite è
condizione necessaria affinché la logica resti testabile su calcolatore e
riusabile al variare delle tecnologie.

\section{Copertura dei requisiti}
La Tabella~\ref{tab:copertura} riassume il legame tra livelli di validazione e
requisiti funzionali. Alcuni requisiti sono coperti dalle prove descritte, altri
restano da validare in modo end-to-end e sono ricondotti agli sviluppi futuri,
in coerenza con la natura di prima iterazione del lavoro.

\begin{table}[h]
\centering
\caption{Copertura dei requisiti funzionali da parte della validazione.}
\label{tab:copertura}
\small
\begin{tabular}{p{2.0cm} p{6.5cm} p{4.0cm}}
\toprule
\textbf{Requisito} & \textbf{Evidenza di validazione} & \textbf{Stato} \\
\midrule
RF-1 & Test host + ingresso pick-up reale & Coperto \\
RF-2 & Test host su interpolazione mappe & Coperto \\
RF-4 & Esecuzione all-fake e mista & Coperto \\
RF-6 & Test host su politiche di sicurezza & Coperto \\
RF-7 & Test host su gestione mappe & Coperto \\
RF-3, RF-5 & Schedulazione accensione/attuatori & Da validare su banco \\
RF-8, RF-9 & Telemetria e comandi & Da validare end-to-end \\
RF-10, RF-11 & Sessioni e OTA & Da validare end-to-end \\
\bottomrule
\end{tabular}
\end{table}

% ============================================================
\chapter{Conclusioni e sviluppi futuri}
\label{ch:conclusioni}

\section{Conclusioni}
Il lavoro ha definito e in parte realizzato una ECU \emph{software-defined} per
motori monocilindrici, intesa come prima iterazione e base architetturale. Il
percorso ha seguito le fasi canoniche di un progetto ingegneristico: definizione
dell'idea e del contesto, analisi dei requisiti, progettazione di
un'architettura indipendente dalla realizzazione, sviluppo di un prototipo su un
sottoinsieme di funzionalità e validazione mirata.

Il risultato principale non è una prestazione motoristica, bensì la
dimostrazione che un'architettura modulare, fondata sulla separazione tra dominio
e infrastruttura e sulla netta distinzione tra compiti critici e non critici,
costituisce una base solida e testabile. L'astrazione delle sorgenti sensoriali,
in particolare, abilita una strategia di validazione a più livelli che rende
sostenibile l'evoluzione del software e riprende, in scala ridotta, il paradigma
\emph{hardware-in-the-loop} dell'industria automobilistica.

\section{Riflessioni critiche}
Alcune scelte vanno lette alla luce del loro orizzonte. L'uso di mappe
monodimensionali è adeguato alla prima iterazione ma limita la fedeltà del
controllo rispetto a mappe tridimensionali. La simulazione di parte dei sensori
ha permesso di progredire senza il cablaggio completo, ma sposta in avanti la
validazione su motore reale. La mediazione del client per la comunicazione con
il server semplifica il nodo di controllo, ma introduce una dipendenza dalla
presenza del client per le funzioni remote. Nessuna di queste scelte compromette
l'architettura: tutte sono sostituibili senza riprogettare il sistema, il che è
precisamente il pregio di una progettazione orientata ai confini astratti.

\section{Sviluppi futuri}
Gli sviluppi futuri individuati sono:
\begin{itemize}[noitemsep]
  \item integrazione del sensore di detonazione con correzione attiva
        dell'anticipo;
  \item controllo della valvola di scarico;
  \item passaggio a mappe tridimensionali (giri~$\times$~farfalla);
  \item completamento della validazione end-to-end di telemetria, sessioni e OTA;
  \item realizzazione di una pipeline di integrazione continua di tipo HITL;
  \item progetto di una scheda dedicata (PCB);
  \item integrazione dei sensori oggi simulati e validazione su motore reale;
  \item evoluzione del \emph{digital twin} da descrittivo a predittivo, con
        modelli termici, di rischio detonazione e di stima del carico, da
        validare contro dati fisici adeguati.
\end{itemize}
Questi sviluppi confermano la lettura del progetto come percorso iterativo, in
cui la presente versione costituisce una fondazione coerente su cui innestare,
in modo incrementale, le funzioni mancanti.

% ============================================================
\begin{thebibliography}{9}

\bibitem{linkecu}
Link Engine Management,
\emph{Why an Aftermarket ECU? Selecting an Aftermarket Engine Management System}.
Disponibile online: \url{https://linkecu.com/new-to-ecus/why-an-aftermarket-ecu/}.

\bibitem{patent_prog}
\emph{Programmable Internal Combustion Engine Controller}, US Patent 6,999,869.
United States Patent and Trademark Office.

\bibitem{freertos_idf}
Espressif Systems,
\emph{FreeRTOS (IDF) --- ESP-IDF Programming Guide}.
Disponibile online:
\url{https://docs.espressif.com/projects/esp-idf/en/stable/esp32/api-reference/system/freertos_idf.html}.

\bibitem{freertos_perf}
\emph{Measuring the Performance of FreeRTOS on ESP32 Multi-Core},
ScienceDirect (IFAC-PapersOnLine).
Disponibile online:
\url{https://www.sciencedirect.com/science/article/abs/pii/S2405896322003639}.

\bibitem{hil_mathworks}
MathWorks,
\emph{What Is Hardware-in-the-Loop (HIL)?}
Disponibile online:
\url{https://www.mathworks.com/discovery/hardware-in-the-loop-hil.html}.

\bibitem{hil_framework}
\emph{A Virtual Testing Framework for Real-Time Validation of Automotive
Software Systems Based on Hardware in the Loop and Fault Injection},
PMC.
Disponibile online: \url{https://pmc.ncbi.nlm.nih.gov/articles/PMC11207294/}.

\bibitem{dt_review}
\emph{Digital Twins for Internal Combustion Engines: A Brief Review},
Journal of Energy, Environment \& Sustainable Engineering.
Disponibile online: \url{https://journal.cbiore.id/index.php/jese/article/view/5}.

\bibitem{dt_diesel}
\emph{A Digital Twin for Diesel Engines: Operator-infused Physics-Informed
Neural Networks with Transfer Learning for Engine Health Monitoring}, arXiv.
Disponibile online: \url{https://arxiv.org/html/2412.11967v2}.

\bibitem{sdv_iea}
International Energy Agency,
\emph{Vehicle Software and Software-Defined Vehicles --- Analysis}.
Disponibile online:
\url{https://www.iea.org/reports/vehicle-software-and-software-defined-vehicles}.

\end{thebibliography}

\end{document}
